% ---- introduction.tex ----
% Balanced reframe. Data sources: research/memory/findings.md, site/src/data/results.json.

Scaling has driven most recent progress in language modelling, yet it has also obscured a more basic question: at small scale, under a fixed data budget, which design choices actually buy generalization, and which only seem to until they are tested with matched controls? Small language models trained under the strict budgets of the BabyLM challenge~\cite{Warstadt2023, Hu2024BabyLM} are an unusually clean setting in which to ask this, because the data is held fixed and the confounds that dominate frontier-scale training are largely absent. What survives in this setting is informative: it constrains the design space for data-efficient pretraining and tells us which inductive biases are worth their complexity.

One long-standing source of inductive bias is linguistic theory. Classical Pāṇinian grammar, refined over more than a millennium~\cite{Matilal1990WordWorld}, encodes compositional structure at two levels: morphologically, through segmentation and inflection, and at the argument level, through the kāraka categories that assign each participant a semantic role (agent, patient, instrument, recipient, source, locus) defined by its relation to the action rather than by surface order~\cite{KiparskyStaal1969, Cardona1974Karaka}. Because kāraka roles are semantic, they are in principle language-independent, which makes them a natural target for a role-aware pretraining signal even in English. The intuition is familiar from the inductive-bias literature: a prior that shrinks the effective hypothesis space should reduce the data needed to reach a given level of competence. Whether this intuition holds for masked-language-model pretraining, when tested rather than assumed, is the empirical question of this paper.

Prabhāsa investigates that question through a pre-registered hypothesis hierarchy. The primary hypothesis, H1, asks whether Pāṇinian structure, realized as Vidyut N-hot morpheme-boundary embeddings~\cite{Prasad2024Vidyut} together with kāraka-aware adaptive masking, measurably improves compositional competence on English syntax (the BLiMP agreement and argument-structure subsets) under matched budgets. We decompose H1 into two mechanistically distinct levers so that they can fail or succeed independently. RQ-A (masking causality) asks whether, at matched mask-rate budget, stratifying the masking distribution by kāraka role changes downstream syntax: that is, whether \emph{which} tokens are masked matters once \emph{how many} are masked is held fixed. RQ-B (auxiliary objective) asks whether a supervised token-level kāraka-role prediction head, trained jointly with MLM and discarded at inference (a computational analogue of śābdabodha, verbal cognition), supplies a signal the masking lever does not. Two further hypotheses frame the wider program: H2 (in scope for follow-up) concerns a Pañcāvayava reasoning scaffold, and H3 (out of scope here) concerns constructive epistemic constraints; neither is evaluated in this paper's results.

Our findings reorder the usual expectation. The clearest, most reproducible effect we observe is not from the Pāṇinian mechanisms but from the pretraining objective, and it is scale-dependent. At the 10M-token budget the choice between pure MLM and a hybrid MLM\,+\,CLM objective is immaterial, but at the 100M-token budget pure MLM outperforms the hybrid by 5.49 BLiMP points. This places our results in direct dialogue with GPT-BERT~\cite{CharpentierSamuel2024}, the BabyLM~2024 winner, whose hybrid masked-plus-causal recipe we find is not uniformly preferable as scale grows. Against this positive result, both Pāṇinian levers return pre-registered null or weak outcomes once their confounds are removed: kāraka-stratified masking is causally inert at matched budget, and the kāraka auxiliary objective produces only a small, non-significant gain that attenuates as seeds accumulate. Figure~\ref{fig:forest} summarizes all pre-registered effects on a common scale.

These outcomes are a feature of the methodology, not a disappointment of it. Earlier single-seed development runs had suggested a $+1.23$-point gain from kāraka masking; the matched-budget comparison shows that apparent gain was an artifact of unequal effective mask rates rather than role structure. This is exactly what pre-registration, matched-budget arms, and multi-seed replication are for. Accordingly, the contributions of this work are methodological as much as empirical. First, we test RQ-A and RQ-B as separately controlled arms, isolating the masking distribution from auxiliary supervision and removing the mask-rate confound that has muddied prior single-arm comparisons. Second, we document a clean, scale-dependent objective effect with effect sizes and confidence intervals, and connect it to the current BabyLM state of the art. Third, every hypothesis test uses at least three seeds, paired bootstrap resampling within pre-registered target sets, explicit outcome thresholds fixed before results were seen, and an adversarial ``Tarka'' memo that states the strongest objection to each finding before closure. Fourth, all code, data, checkpoints, and per-seed results are released openly with full configuration hashing.

To our knowledge this is the first pre-registered, multi-seed, mechanistically decomposed test of Pāṇinian structure as an inductive bias for masked-language-model pretraining under strict data budgets. Prior linguistically motivated pretraining, from morphology-aware translation models~\cite{Belinkov2017Morphology} to morphologically grounded word representations~\cite{Botha2014Morphology}, has typically tested a single prior in isolation or at larger scale. The distinctive contribution here is the experimental design, mechanistic separation under matched budgets with pre-registered inference, which yields a trustworthy map of what helps small-model pretraining (objective and architecture) and what does not at this scale (the kāraka masking distribution). A well-documented null that overturns an optimistic prior is a result the field needs more of, and we report ours in full. The remainder of the paper is organized as follows. Section~\ref{sec:related} reviews data-efficient pretraining, morphology-aware masking, and computational treatments of the Pāṇinian and Navya-Nyāya traditions. Section~\ref{sec:data} describes the corpora and the four pre-pretraining dose arms. Section~\ref{sec:methods} presents the architecture, the three Pāṇinian mechanisms, and the pre-registered research questions. Section~\ref{sec:training} details the training and evaluation protocol. Section~\ref{sec:results} reports each finding with effect estimates and intervals. Sections~\ref{sec:discussion} and~\ref{sec:conclusion} interpret the results, state threats to validity, and outline future work.
